% =========================================================================
\section{Introducción: qué estudiamos en la sesión}\label{sec:introduccion}
% =========================================================================
Un \emph{amplificador operacional} (AO) es un amplificador de tensión diferencial con una
ganancia tan grande que, por sí solo, apenas es útil: cualquier diferencia de milivoltios
entre sus entradas lo lleva a saturación. Su interés está en la \emph{realimentación
negativa}: si se devuelve una fracción de la salida a la entrada inversora mediante una red
pasiva, el comportamiento del conjunto deja de depender del AO y pasa a depender solo de esa
red. La sesión~1 del laboratorio~\cite{GUION-S1} explota esta idea en tres circuitos
(\cref{fig:esquema}):
\begin{enumerate}
  \item un \emph{amplificador inversor} de ganancia $A_V=-10$;
  \item un \emph{amplificador no inversor} de ganancia $A_V=+11$;
  \item un \emph{amplificador inversor generalizado}, en el que las resistencias se sustituyen
        por impedancias. Si la red de realimentación contiene un condensador, la ganancia
        depende de la frecuencia y el circuito es un \emph{filtro}; en el dominio del tiempo
        integra o deriva la señal, y la transformada de Laplace enlaza ambas descripciones.
\end{enumerate}
La parte~I reúne las herramientas imprescindibles, con la notación de la teoría de la
asignatura~\cite{TEORIA-EA,TEORIA-ADICIONAL,TEORIA-LAPLACE} y del texto de referencia
de Sedra y Smith~\cite{SEDRA}; la parte~II resuelve los recuadros del PreLab.

% Figura: estructura lógica de la sesión.
\begin{figure}[ht]
\centering
\begin{tikzpicture}[
    caja/.style={draw=azulclaro, fill=azulfondo, line width=0.5pt, align=center,
                 inner sep=5pt, font=\small, minimum height=1.0cm},
    final/.style={draw=azulacero, fill=white, line width=0.8pt, align=center,
                  inner sep=5pt, font=\small, minimum height=1.0cm},
    fl/.style={->, line width=0.7pt, azulacero, shorten >=2pt, shorten <=2pt},
    nota/.style={font=\scriptsize, text=grispizarra, align=center}]
  \node[caja] (M) at (0,0) {Modelo ideal del AO\\[-1pt]{\footnotesize $i_\pm=0,\ v_+=v_-$}};
  \node[caja] (L) at (0,-1.9) {Impedancias\\[-1pt]{\footnotesize $Z_R=R,\ Z_C=1/(sC)$}};
  \node[caja] (G) at (4.4,-0.95) {Inversor generalizado\\[-1pt]{\footnotesize $H(s)=-Z_2/Z_1$}};
  \node[final] (I) at (9.3,-0.95) {Inversor\\[-1pt]{\footnotesize $A_V=-R_2/R_1$}};
  \node[final] (N) at (9.3,0.85) {No inversor\\[-1pt]{\footnotesize $A_V=1+R_2/R_1$}};
  \node[final] (F) at (9.3,-2.65) {Filtro paso bajo\\[-1pt]{\footnotesize $Z_2=R_F\parallel C$}};
  \draw[fl] (M.east) -- (G.north west);
  \draw[fl] (L.east) -- (G.south west);
  \draw[fl] (G.east) -- (I.west);
  \draw[fl] (G.east) -- (F.west);
  \draw[fl] (M.east) ++(0,0.12) -- ($(N.west)+(0,0)$);
  \node[nota] at (6.95,-0.72) {$Z_k=R_k$};
  \node[nota] at (6.75,-2.25) {$Z_1=R$};
\end{tikzpicture}
\caption{Estructura lógica de la sesión. El modelo ideal del AO, junto con el concepto de
impedancia, da la ganancia del amplificador inversor generalizado; el inversor y el filtro
paso bajo son casos particulares. El no inversor se resuelve con las mismas reglas.}
\label{fig:esquema}
\end{figure}

